\documentclass[11pt]{article}

\usepackage{amsmath,amssymb,amsthm,mathtools}
\usepackage[margin=1in]{geometry}
\usepackage{hyperref}
\usepackage{enumitem}
\usepackage{cleveref}
\usepackage{graphicx}

\newtheorem{theorem}{Theorem}[section]
\newtheorem{proposition}[theorem]{Proposition}
\newtheorem{lemma}[theorem]{Lemma}
\newtheorem{corollary}[theorem]{Corollary}
\theoremstyle{definition}
\newtheorem{definition}[theorem]{Definition}
\newtheorem{example}[theorem]{Example}
\newtheorem{assumption}[theorem]{Assumption}
\theoremstyle{remark}
\newtheorem{remark}[theorem]{Remark}

\DeclareMathOperator{\Tr}{Tr}
\DeclareMathOperator{\im}{im}
\newcommand{\E}{\mathbb{E}}
\newcommand{\R}{\mathbb{R}}
\newcommand{\C}{\mathbb{C}}
\newcommand{\N}{\mathbb{N}}
\newcommand{\HH}{\mathcal{H}}
\newcommand{\FF}{\mathcal{F}}
\newcommand{\DD}{\mathcal{D}}
\newcommand{\ip}[2]{\langle #1, #2 \rangle}
\newcommand{\norm}[1]{\| #1 \|}

\title{The Fractional Memory Process:\\
Well-Posedness, Spectral Decomposition,\\
and Discrete Representation}
\author{Ramiro Fontes\\Quijotic Research LLC\\{\small\texttt{ramiro@quijotic.com}}}
\date{\today}

\begin{document}
\maketitle

\begin{abstract}
We study the fractional stochastic evolution equation
\[
\partial_t^{2H_2} X_t = A_{H_1} X_t + \dot{B}_t,
\qquad X_0 = f,
\]
where $\partial_t^{2H_2}$ is the Caputo fractional
derivative, $A_{H_1} = -c_{H_1}(-\Delta)^{H_1}$ is
the fractional Laplacian ($c_{H_1} > 0$ a constant),
$\dot{B}_t$ is Gaussian noise, and $f \in L^2(\R)$
is the initial datum.  The two parameters
$H_1 \in (0,1]$ and $H_2 \in (1/4,1)$ independently
control spatial resolution and temporal persistence.

We establish well-posedness of the mild formulation
under trace-class noise for all admissible parameters
and under cylindrical noise when additionally
$H_1 > \alpha/(2(2\alpha-1))$ with $\alpha = 2H_2$
(which forces $H_2 > 1/3$).  The spectral decomposition reduces the
SPDE to independent scalar fractional
Ornstein--Uhlenbeck modes with power-law decay
$t^{-H_2/(2H_1)}$ (sharp, with matching lower bound
when $\hat{f}(0) \neq 0$), and the variance scales
with explicit self-similarity exponent
$H = 2H_2 - 1/2 - H_2/(2H_1)$
(not to be confused with the parameters $H_1$, $H_2$).

For a stream of discrete impulse inputs, we prove
that the Fourier-projected mild solution is a
finite sum of two-parameter Mittag-Leffler
evaluations with explicit prefactors.  A distinct
autonomous-decay representation using the
one-parameter Mittag-Leffler function gives the
free resolvent of the fractional Ornstein--Uhlenbeck
modes.  The two representations coincide at
$\alpha = 1$ and play complementary roles for
$\alpha \in (0, 1)$: the forced-SPDE Green's
function is $L^1$-integrable on $[0, \infty)$ with
total mass $1/\lambda$ but has a $(t-s)^{\alpha-1}$
singularity at the impulse time, while the
autonomous-decay kernel is regular at zero age but
$L^1$-divergent at infinity.  Under stationary
write processes this produces a power-law
$T^{1-\alpha}$ separation in aggregate kernel mass,
established in closed form; in a retrieval pipeline the
singularity causes the forced-SPDE kernel to
dominate semantic signals, also verified
numerically.  The effective temporal depth under
modal truncation to $N$ modes scales as
$N^{H_1/H_2}$.

We also identify the operator derivative structure
on the fractional memory space, yielding exact
Gaussian representations and a polynomially
propagating Leibniz defect under L\'evy noise.
These results furnish the kernel and scaling laws
used in the companion architecture
paper~\cite{fontes2026arch}.
\end{abstract}

\section{Introduction}\label{sec:intro}

The stochastic heat equation
$\partial_t X = \Delta X + \dot{B}$ governs
diffusion-driven dynamics in infinite dimensions and
is the foundational object of Da~Prato--Zabczyk
theory~\cite{daprato2014}.  Its finite-dimensional
projections are classical Ornstein--Uhlenbeck
processes: Markov, Gaussian, with exponential
relaxation to equilibrium.  These properties make
the theory analytically tractable but impose a
structural limitation: the Markov property encodes
history only through the current state, and
relaxation is exponentially fast.  This is a
structural mismatch for systems whose function
depends on retaining information across long
timescales --- viscoelastic materials, anomalous
transport, and the memory-architecture applications
motivating the companion paper~\cite{fontes2026arch},
in which the polynomial relaxation and exact
discrete-impulse representation derived below
furnish the kernel and scaling laws used at
runtime.

A growing body of work---in anomalous diffusion,
viscoelasticity, and fractional
kinetics~\cite{mainardi2010,meerschaert2012}---replaces
the integer-order time derivative with a Caputo
fractional derivative $\partial_t^\alpha$, producing
power-law relaxation via the Mittag-Leffler function
$E_\alpha(-\lambda t^\alpha)$.  Independently, the
integer-order Laplacian is replaced by a fractional
power $(-\Delta)^s$ in models of anomalous
spatial transport.

This paper studies the equation that combines both
generalizations:
\begin{equation}\label{eq:quijotic_main}
\partial_t^{2H_2} X_t = A_{H_1} X_t + \dot{B}_t,
\qquad X_0 = f,
\end{equation}
where $\partial_t^{2H_2}$ is the Caputo derivative
with $\alpha = 2H_2 \in (1/2, 2)$,
$A_{H_1} = -c_{H_1}(-\Delta)^{H_1}$ is the
fractional Laplacian with $H_1 \in (0,1]$, and
$\dot{B}_t$ is Gaussian noise on $L^2(\R)$.
The two parameters are independent: $H_1$ controls
the spatial coupling range and $H_2$ controls
temporal persistence.

\subsection{Main results}

\begin{enumerate}[label=\textbf{(\Alph*)}]
\item \textbf{Well-posedness}
  (\Cref{thm:existence_trace,thm:existence}).
  Unique mild solutions exist for all
  $H_1 \in (0,1]$ and $H_2 \in (1/4,1)$ with
  trace-class noise; cylindrical noise requires
  additionally $H_1 > \alpha/(2(2\alpha-1))$ with
  $\alpha = 2H_2$ (which forces $H_2 > 1/3$).

\item \textbf{Spectral decomposition}
  (\Cref{thm:spectral}).  The SPDE reduces to
  independent scalar fractional OU processes
  with mode-dependent decay rate
  $\lambda(\xi) = c_{H_1}|\xi|^{2H_1}$
  (where $\xi$ is the Fourier variable).

\item \textbf{Energy and decay}
  (\Cref{thm:energy,thm:decay_det}).  The
  deterministic component decays as
  $t^{-H_2/(2H_1)}$---polynomial for all
  parameters, with matching lower bound.

\item \textbf{Variance scaling}
  (\Cref{thm:selfsimilar}).  Explicit scaling
  exponent
  $H = 2H_2 - 1/2 - H_2/(2H_1)$.

\item \textbf{Discrete-impulse exactness}
  (\Cref{thm:discretization}).  For impulse
  inputs, the Fourier-projected mild solution
  is a finite Mittag-Leffler sum, and modal
  truncation to $N$ modes gives effective
  temporal depth $\sim N^{H_1/H_2}$
  (\Cref{thm:eff_context}).

\item \textbf{Operator derivative and defect
  structure}
  (\Cref{thm:clark_ocone,thm:defect_dichotomy}).
  The operator derivative $D = \delta^*$ yields
  exact Clark--Ocone representation in the Gaussian
  case; under L\'evy noise, the Leibniz defect is
  nonzero, concentrated on the jump component, and
  propagates polynomially via the Mittag-Leffler
  resolvent kernel.
\end{enumerate}

\subsection{Related work}

Fractional stochastic evolution equations have been
studied by Mainardi~\cite{mainardi2010} (fractional
relaxation), Meerschaert--Sikorskii~\cite{meerschaert2012}
(stochastic fractional PDEs), and
Da~Prato--Zabczyk~\cite{daprato2014}
(infinite-dimensional stochastic equations).  The
present work differs in three respects: combining
fractional time derivatives with fractional spatial
operators under independent parameters; connecting
the resulting dynamics to an operator calculus
(representation and defect structure); and proving the
exactness of the discrete-impulse representation that
motivates computational applications.

The operator derivative framework
$D = \delta^*$~\cite{fontes2026opfac,fontes2026banach}
and its application to Volterra-L\'evy
markets~\cite{fontes2026volterra} provide the calculus
used in Sections~\ref{sec:operator}--\ref{sec:levy}
to derive the defect dichotomy and its propagation law
for the fractional memory process.

\section{Preliminaries}\label{sec:prelim}

We work on a filtered probability space
$(\Omega, \FF, (\FF_t)_{t \geq 0}, \mathbb{P})$
satisfying the usual conditions, with noise driven by
either a $Q$-Wiener process (trace-class covariance)
or a cylindrical Brownian motion on $L^2(\R)$.
A \emph{cylindrical Brownian motion} is a family
$(B_t)_{t \geq 0}$ with
$\E[(B_t h)^2] = t\norm{h}^2$ for all
$h \in L^2(\R)$; a \emph{$Q$-Wiener process}
replaces the identity covariance with a trace-class
operator $Q$, giving
$\E[(B_t^Q h_1)(B_s^Q h_2)]
= (t \wedge s)\ip{Qh_1}{h_2}$.
For an operator-valued integrand
$\Psi_s : L^2(\R) \to L^2(\R)$, the Hilbert-Schmidt
norm
$\norm{\Psi}_{\mathrm{HS}}^2
:= \sum_k \norm{\Psi e_k}_{L^2}^2$
(summed over any orthonormal basis $(e_k)$ of
$L^2(\R)$) controls integrability of the stochastic
convolution:
$\int_0^t \Psi_s\, dB_s$ is well-defined when
$\int_0^t \norm{\Psi_s}_{\mathrm{HS}}^2\, ds < \infty$
a.s.

The time derivative in~\eqref{eq:quijotic_main} is
the Caputo fractional derivative.  For
$0 < \alpha < 1$ and $u \in AC([0,T])$,
\begin{equation}\label{eq:caputo}
\partial_t^\alpha u(t)
:= \frac{1}{\Gamma(1-\alpha)}
\int_0^t (t-s)^{-\alpha} u'(s)\, ds;
\end{equation}
for $1 < \alpha < 2$ and $u \in AC^2([0,T])$,
the first derivative is replaced by the second.
(When $\alpha = 1$, $\partial_t^\alpha = \partial_t$.)
Since $\alpha = 2H_2$ with $H_2 \in (1/4,1)$,
both regimes arise.
Unlike the ordinary derivative,
$\partial_t^\alpha(uv) \neq
u\,\partial_t^\alpha v + v\,\partial_t^\alpha u$
in general; this failure of the Leibniz rule
reappears in the energy analysis of
\Cref{sec:energy}.  The Caputo derivative also
satisfies $\partial_t^\alpha c = 0$ for constants.
See~\cite{podlubny1999} for a comprehensive
treatment.

The resolvent of the fractional equation is built
from the Mittag-Leffler function
$E_{\alpha,\beta}(z)
:= \sum_{k=0}^\infty z^k/\Gamma(\alpha k + \beta)$,
with $E_\alpha := E_{\alpha,1}$.  Three properties
are used throughout:
$E_\alpha(-t^\alpha)$ is completely monotone for
$\alpha \in (0,1)$~\cite{pollard1948};
$E_1(z) = e^z$ (recovering the exponential);
and $E_\alpha(-r) \sim r^{-1}/\Gamma(1-\alpha)$
as $r \to \infty$~\cite{gorenflo2014}, which gives
the power-law tail that distinguishes the fractional
case from the Markov one.  We also use the Laplace
identity
$\mathcal{L}[t^{\beta-1}
E_{\alpha,\beta}(\pm\lambda t^\alpha)](s)
= s^{\alpha-\beta}/(s^\alpha \mp \lambda)$.

The following bounds across the full parameter range
$\alpha \in (0, 2)$ are used repeatedly in the proofs
that follow; we collect them here so that the
arguments below do not need to treat $\alpha \leq 1$
and $\alpha > 1$ separately.

\begin{lemma}[Uniform Mittag-Leffler bounds]
\label{lem:ml_bounds}
For every $\alpha \in (0, 2]$ there exists a
constant $C_\alpha > 0$ such that for all $r \geq 0$,
\begin{enumerate}[label=\textnormal{(\roman*)}]
\item $|E_\alpha(-r)| \leq 1$;
\item $|E_\alpha(-r)| \leq C_\alpha/(1 + r)$;
\item $|E_{\alpha,\alpha}(-r)| \leq C_\alpha/(1 + r)$.
\end{enumerate}
\end{lemma}

\begin{proof}
For $\alpha \in (0, 1]$, (i) is complete monotonicity
of $E_\alpha(-t^\alpha)$~\cite{pollard1948,schneider1996},
which gives the stronger statement
$0 \leq E_\alpha(-r) \leq 1$.  For $\alpha \in (1, 2]$
the function $E_\alpha(-r)$ is not monotone and takes
negative values, but the bound $|E_\alpha(-r)| \leq 1$
still holds via the integral representation of
$E_\alpha$ as a Laplace transform of a finite measure
of total mass one; see \cite[Prop.~3.10, \S4.3]{gorenflo2014}.
The boundary case is
$E_2(-r) = \cos(\sqrt{r}) \in [-1, 1]$.

Parts (ii) and (iii) follow from the Mittag-Leffler
asymptotic~\cite[Ch.~4]{gorenflo2014}
\[
E_{\alpha,\beta}(-r) = -\sum_{k=1}^{N}
\frac{(-r)^{-k}}{\Gamma(\beta - \alpha k)}
+ O(r^{-N-1}),
\qquad r \to \infty, \quad \alpha \in (0, 2),
\]
combined with continuous-function boundedness on
$[0, R]$ for any $R > 0$.  For (iii) the leading
asymptotic coefficient vanishes (since
$\Gamma(\alpha - \alpha)^{-1} = 0$), so in fact
$E_{\alpha,\alpha}(-r) = O(r^{-2})$; we use only
the weaker $C_\alpha/(1+r)$ bound below.
\end{proof}

\begin{remark}[Sharp asymptotic for
$E_{\alpha,\alpha}$]
For $\alpha \in (0, 1)$, the relation
$(t)^{\alpha-1} E_{\alpha,\alpha}(-\lambda t^\alpha)
= -\lambda^{-1}\,\partial_t
E_\alpha(-\lambda t^\alpha)$
combined with the $E_\alpha$ asymptotic gives the
sharper bound
$E_{\alpha,\alpha}(-r) \sim \alpha/
[\Gamma(1-\alpha)\,r^2]$
as $r \to \infty$.  We use only the weaker
$C_\alpha/(1+r)$ bound in the proofs below.
\end{remark}

The spatial operator
$A_{H_1} = -c_{H_1}(-\Delta)^{H_1}$, where
$c_{H_1} > 0$ is a fixed constant, acts in Fourier
space (we write $\hat{u}$ for the Fourier transform
of $u$) as
$\widehat{(-\Delta)^{H_1} u}(\xi)
= |\xi|^{2H_1} \hat{u}(\xi)$ for
$H_1 \in (0,1]$, and generates an analytic contraction
semigroup with Fourier symbol
$e^{-c_{H_1}|\xi|^{2H_1} t}$.

\section{Well-Posedness}\label{sec:wellposed}

We formulate the mild solution
of~\eqref{eq:quijotic_main} via a resolvent family
in Fourier space.  Write $\alpha = 2H_2$ throughout
this section.

\begin{definition}[Resolvent family]\label{def:resolvent}
$\widehat{S_t^{(\alpha, H_1)} f}(\xi) = E_\alpha(-c_{H_1}|\xi|^{2H_1} t^\alpha)\,\hat{f}(\xi)$.
\end{definition}

\begin{lemma}[Resolvent properties]\label{lem:resolvent_props}
(a) $S_0 = I$. (b) $\norm{S_t f}_{L^2} \leq \norm{f}_{L^2}$. (c) $t \mapsto S_t f$ continuous. (d) $\partial_t^\alpha S_t f = A_{H_1} S_t f$ for $f \in \DD(A_{H_1}) := H^{2H_1}(\R)$.
\end{lemma}

\begin{proof}
(a) $E_\alpha(0) = 1$. (b) By \Cref{lem:ml_bounds}(i), $|E_\alpha(-r)| \leq 1$; apply Plancherel. (c) Dominated convergence. (d) Mittag-Leffler differentiation.
\end{proof}

The stochastic convolution requires a second kernel.

\begin{definition}[Resolvent integrand]\label{def:resolvent_kernel}
$\widehat{\Phi_t}(\xi) := t^{\alpha-1} E_{\alpha,\alpha}(-c_{H_1}|\xi|^{2H_1} t^\alpha)$.
\end{definition}

The mild solution is the sum of the deterministic
resolvent applied to the initial datum and the
stochastic convolution against the noise:

\begin{definition}[Mild solution]\label{def:mild}
$X_t = S_t^{(\alpha, H_1)} f + \int_0^t \Phi_{t-s}\, dB_s$.
\end{definition}

\subsection{Existence and uniqueness}

\begin{theorem}[Trace-class noise]
\label{thm:existence_trace}
For $f \in L^2(\R)$, $H_1 \in (0,1]$, $H_2 \in (1/4,1)$,
$T > 0$, and $Q$-Wiener noise with trace-class $Q$,
equation~\eqref{eq:quijotic_main} admits a unique mild
solution $X \in L^2(\Omega; C([0,T]; L^2(\R)))$.
\end{theorem}

\begin{proof}
\textbf{Step~1.} $S_t f \in C([0,T]; L^2)$ by
\Cref{lem:resolvent_props}.

\textbf{Step~2 (Hilbert-Schmidt condition).}
The resolvent integrand satisfies
$\norm{\Phi_s}_{L^2 \to L^2}
= \sup_\xi |\widehat{\Phi_s}(\xi)|
\leq C s^{\alpha-1}$
(from $|E_{\alpha,\alpha}(-r)| \leq C/(1+r)$,
evaluated at $r = 0$).  Therefore
\[
\norm{\Phi_s Q^{1/2}}_{\mathrm{HS}}^2
\leq \norm{\Phi_s}_{L^2 \to L^2}^2 \Tr(Q)
\leq C^2 s^{2(\alpha-1)}\,\Tr(Q).
\]
The time integral converges:
$\int_0^t s^{2(\alpha-1)}\, ds < \infty$
when $2(\alpha-1) > -1$, i.e., $\alpha > 1/2$, i.e.,
$H_2 > 1/4$.

\textbf{Step~3.} Continuity: standard stochastic
convolution estimates~\cite{daprato2014}.

\textbf{Step~4 (Uniqueness).}  The mild-solution
formula in \Cref{def:mild} is a specific functional
of $f$ and $B$; any two mild solutions with the same
data are equal by construction.  The map
$(f, B) \mapsto X$ is linear, so the difference of
two solutions with common data satisfies the
homogeneous equation with zero initial datum and
zero forcing, which is identically zero.
\end{proof}

\begin{remark}[The condition $H_2 > 1/4$]
\label{rem:H2_quarter}
The restriction $H_2 > 1/4$ arises from the temporal
singularity $s^{\alpha-1}$ of the resolvent integrand
at $s = 0$.  For $H_2 \leq 1/4$, existence can be
recovered by imposing additional spatial regularity on
the noise covariance (e.g., requiring $Q : L^2 \to H^s$
for $s$ large enough that the Mittag-Leffler spatial
decay compensates the temporal singularity).  We do not
pursue this here; the range $H_2 \in (1/4, 1)$ is
sufficient for the results below.
\end{remark}

\begin{corollary}[Cylindrical noise]
\label{thm:existence}
Let $\alpha = 2H_2$.  For $f \in L^2(\R)$,
$H_2 \in (1/4,1)$, and
$H_1 \in (\alpha/(2(2\alpha-1)),\ 1]$
(which forces $\alpha > 2/3$, equivalently
$H_2 > 1/3$), equation~\eqref{eq:quijotic_main}
with cylindrical Brownian noise ($Q = I$) admits a
unique mild solution.
\end{corollary}

\begin{proof}
With cylindrical noise, $\norm{\Phi_s Q^{1/2}}_{\mathrm{HS}}
= \norm{\Phi_s}_{\mathrm{HS}}$.  Using
$|E_{\alpha,\alpha}(-r)| \leq C/(1+r)$:
\[
\norm{\Phi_s}_{\mathrm{HS}}^2
\leq C\,s^{2(\alpha-1)} \int_\R
\frac{d\xi}{(1 + c|\xi|^{2H_1} s^\alpha)^2}.
\]
Substituting $\eta = |\xi|^{2H_1} s^\alpha$:
the spatial integral is
$\leq C\,s^{-\alpha/(2H_1)}
\int_0^\infty \eta^{1/(2H_1)-1}(1+c\eta)^{-2}\, d\eta$,
which is the spatial-integrability requirement and
converges when $H_1 > 1/4$.  Thus
$\norm{\Phi_s}_{\mathrm{HS}}^2
\leq C\,s^{2\alpha - 2 - \alpha/(2H_1)}$, and time
integrability $\int_0^t \norm{\Phi_s}_{\mathrm{HS}}^2\,
ds < \infty$ near $s = 0$ requires this exponent to
exceed $-1$, i.e.\ $2\alpha - 2 - \alpha/(2H_1) > -1$,
equivalently $H_1 > \alpha/(2(2\alpha - 1))$.  Since
$\alpha/(2(2\alpha - 1)) > 1/4$ for every
$\alpha > 1/2$, this stochastic-convolution condition
is strictly stronger than the spatial-integral
condition $H_1 > 1/4$ used above and supersedes it.
The combined hypothesis is therefore
\[
H_1 > \max\bigl\{1/4,\ \alpha/(2(2\alpha-1))\bigr\}
\;=\; \alpha/(2(2\alpha-1)),
\]
which together with $H_1 \leq 1$ requires
$\alpha > 2/3$, i.e.\ $H_2 > 1/3$.
\end{proof}

\subsection{Pathwise continuity}

\begin{theorem}\label{thm:continuity}
Under the assumptions of \Cref{thm:existence_trace}
or \Cref{thm:existence}, the mild solution satisfies
$X \in C([0,T]; L^2(\R))$ almost surely.
\end{theorem}

\begin{proof}
The deterministic part is continuous by
\Cref{lem:resolvent_props}(c).  For the stochastic
convolution, we use the factorization
method~\cite{daprato2014}.

Choose $\gamma \in (0, 1/2)$ with
$\gamma < \alpha - 1/2$; this requires
$\alpha > 1/2$, i.e., $H_2 > 1/4$.  (Under trace-class noise, the $Q$-weighted norm
provides the same integrability within the range
$H_2 > 1/4$ established in \Cref{thm:existence_trace}.)

Define $Y_t := \int_0^t (t-s)^{-\gamma}
\Phi_{t-s}\, dB_s$.  By the Burkholder-Davis-Gundy
inequality,
$\E[\norm{Y_t}^{2p}]
\leq C_p \left(\int_0^t s^{-2\gamma}
\norm{\Phi_s}_{\mathrm{HS}}^2\, ds\right)^p < \infty$
under the integrability conditions of the applicable
existence theorem.

The stochastic convolution is recovered via
$\int_0^t \Phi_{t-s}\, dB_s
= \frac{\sin(\pi\gamma)}{\pi}
\int_0^t (t-s)^{\gamma-1} Y_s\, ds$.
The fractional integral maps $L^\infty$ to
$C^{\gamma - 1/p}$, giving a.s.\ H\"older continuity.
\end{proof}

\section{Spectral Representation}\label{sec:spectral}

\begin{theorem}[Spectral decomposition]\label{thm:spectral}
The mild solution has Fourier representation
\begin{equation}\label{eq:spectral}
\hat{X}_t(\xi) = E_\alpha(-\lambda(\xi)\,t^\alpha)\,\hat{f}(\xi) + \int_0^t (t-s)^{\alpha-1} E_{\alpha,\alpha}(-\lambda(\xi)\,(t-s)^\alpha)\, d\hat{B}_s(\xi),
\end{equation}
where $\lambda(\xi) = c_{H_1}|\xi|^{2H_1}$ is the
mode-dependent decay rate and
$\hat{B}_t(\xi) := \int_\R e^{-i\xi x} dB_t(x)$
is the Fourier-space Brownian motion, satisfying
$\E[d\hat{B}_t(\xi)\,\overline{d\hat{B}_s(\eta)}]
= \delta(\xi - \eta)\,\delta(t-s)\,d\xi\,dt$.
\end{theorem}

\begin{proof}
Fourier transform of the mild solution; the resolvent and stochastic integral are Fourier multipliers.
\end{proof}

\begin{corollary}[Mode-wise dynamics]\label{cor:modewise}
Each mode evolves as a scalar fractional OU: $\partial_t^\alpha \hat{X}_t(\xi) = -\lambda(\xi)\hat{X}_t(\xi) + d\hat{B}_t(\xi)/dt$. High-frequency modes decay rapidly; low-frequency modes persist.
\end{corollary}

\begin{corollary}[Mode-wise variance]\label{cor:mode_variance}
For $f = 0$: $\E[|\hat{X}_t(\xi)|^2] = \int_0^t s^{2(\alpha-1)} |E_{\alpha,\alpha}(-\lambda(\xi)s^\alpha)|^2\, ds$.
\end{corollary}

\section{Regularity}\label{sec:regularity}

\begin{theorem}[Spatial regularity]\label{thm:spatial_reg}
Let $\alpha = 2H_2$.
\begin{enumerate}[label=\textnormal{(\alph*)}]
\item \emph{Deterministic smoothing.}  For $f \in L^2(\R)$
and $t > 0$, $S_t f \in H^s(\R)$ for every
$0 \leq s \leq 2H_1$, with
$\norm{S_t f}_{H^s} \leq C(s,t,H_1,H_2)\,\norm{f}_{L^2}$.
\item \emph{Trace-class noise.}  If $Q$ is trace-class,
the stochastic convolution
$\int_0^t \Phi_{t-r}\, dB_r$ belongs to $H^s(\R)$ provided
\begin{equation}\label{eq:spatial_reg_trace}
\int_0^t \norm{(I-\Delta)^{s/2} \Phi_r Q^{1/2}}_{\mathrm{HS}}^2\,
dr < \infty.
\end{equation}
Without additional smoothing of $Q$, the operator-norm
estimate
$\norm{(I-\Delta)^{s/2}\Phi_r Q^{1/2}}_{\mathrm{HS}}
\leq \norm{(I-\Delta)^{s/2}\Phi_r}_{\mathrm{op}}
\Tr(Q)^{1/2}$
gives the sufficient condition
$s < H_1(2\alpha-1)/\alpha$.
If $(I-\Delta)^{s/2} Q^{1/2}$ is itself Hilbert-Schmidt,
then $H^s$ regularity follows for all $\alpha > 1/2$.
\item \emph{Cylindrical noise.}  For $Q = I$, the
stochastic convolution belongs to $H^s(\R)$ only under the
corresponding Hilbert-Schmidt condition
\begin{equation}\label{eq:spatial_reg_cyl}
\int_0^t \int_\R (1+|\xi|^2)^s\, r^{2\alpha-2}\,
|E_{\alpha,\alpha}(-c_{H_1}|\xi|^{2H_1} r^\alpha)|^2\,
d\xi\, dr < \infty.
\end{equation}
Using the same substitution as in
\Cref{thm:existence}, this is guaranteed when
$s < H_1(2\alpha-1)/\alpha - 1/2$.  At $s = 0$ this
recovers the cylindrical-noise well-posedness threshold
$H_1 > \alpha/(2(2\alpha-1))$.
\end{enumerate}
\end{theorem}

\begin{proof}
(a) On the Fourier side,
$\norm{S_t f}_{H^s}^2 = \int_\R (1+|\xi|^2)^s
|E_\alpha(-c_{H_1}|\xi|^{2H_1} t^\alpha)|^2 |\hat f(\xi)|^2\,
d\xi$.  For $|\xi| \leq 1$, $(1+|\xi|^2)^s \leq 2^s$ and
$|E_\alpha| \leq 1$.  For $|\xi| > 1$,
$(1+|\xi|^2)^s \leq 2^s |\xi|^{2s}$ and
$|E_\alpha(-c_{H_1}|\xi|^{2H_1} t^\alpha)|^2 \leq C\,
|\xi|^{-4H_1} t^{-2\alpha}$, so the high-frequency
contribution is finite when $2s - 4H_1 \leq 0$.

(b) The Hilbert-Schmidt criterion~\eqref{eq:spatial_reg_trace}
is the standard isometry condition for the stochastic
convolution.  Since $\norm{T A}_{\mathrm{HS}}
\leq \norm{T}_{\mathrm{op}}\norm{A}_{\mathrm{HS}}$ and
$\norm{Q^{1/2}}_{\mathrm{HS}}^2 = \Tr(Q)$, the
operator-norm route reduces to bounding
$\norm{(I-\Delta)^{s/2}\Phi_r}_{\mathrm{op}}
= \sup_\xi (1+|\xi|^2)^{s/2}|\widehat{\Phi_r}(\xi)|$.
For $s < 2H_1$, the sup is attained at
$|\xi_*| \sim r^{-\alpha/(2H_1)}$, giving
$\norm{(I-\Delta)^{s/2}\Phi_r}_{\mathrm{op}}^2
\leq C\, r^{2\alpha - 2 - \alpha s/H_1}$;
near $r = 0$ time integrability requires
$2\alpha - 1 - \alpha s/H_1 > 0$, i.e.\
$s < H_1(2\alpha - 1)/\alpha$.  If
$(I-\Delta)^{s/2}Q^{1/2}$ is HS, instead use
$\norm{(I-\Delta)^{s/2}\Phi_r Q^{1/2}}_{\mathrm{HS}}
\leq \norm{\Phi_r}_{\mathrm{op}}
\norm{(I-\Delta)^{s/2}Q^{1/2}}_{\mathrm{HS}}
\leq C\,r^{\alpha-1}\norm{(I-\Delta)^{s/2}Q^{1/2}}_{\mathrm{HS}}$,
and time integrability $\int_0^t r^{2(\alpha-1)}\,dr
< \infty$ holds for all $\alpha > 1/2$.

(c) Apply the substitution
$\eta = c_{H_1}|\xi|^{2H_1} r^\alpha$ to the inner
integrand of~\eqref{eq:spatial_reg_cyl}; the spatial
$\xi$-integral evaluates to
$C\, r^{-\alpha(1+2s)/(2H_1)}$ times a finite
$\eta$-integral (convergent at infinity for
$s < 2H_1 - 1/2$, automatic in the relevant range, and
trivially at $\eta = 0$).  Thus the inner integrand is
bounded by $C\,r^{2\alpha - 2 - \alpha(1+2s)/(2H_1)}$, and
time integrability at $r = 0$ requires
$2\alpha - 1 - \alpha(1+2s)/(2H_1) > 0$, equivalently
$s < H_1(2\alpha - 1)/\alpha - 1/2$.  At $s = 0$ this is
$H_1 > \alpha/(2(2\alpha-1))$, the threshold of
\Cref{thm:existence}.
\end{proof}

\begin{theorem}[Temporal H\"older continuity]\label{thm:temporal_reg}
$\E[\norm{X_t - X_s}_{L^2}^2] \leq C|t-s|^{2\beta}$ with $\beta = \min(2H_2, 1/2 - \varepsilon)$.
\end{theorem}

\begin{proof}
\emph{Deterministic part.}  On the Fourier side,
$\widehat{S_t f}(\xi) - \widehat{S_s f}(\xi)
= [E_\alpha(-\lambda(\xi) t^\alpha)
- E_\alpha(-\lambda(\xi) s^\alpha)]\,\hat{f}(\xi)$.
The Mittag-Leffler function $E_\alpha$ is
differentiable on $(-\infty, 0]$ with
$|E_\alpha'(-r)| \leq C$ on any compact set of $r$
(from the series) and $|E_\alpha'(-r)| = O(r^{-2})$
at infinity (from \Cref{lem:ml_bounds}(iii) and
the identity
$E_\alpha'(-r) = -E_{\alpha,\alpha}(-r)/\alpha$).
So $|E_\alpha(-\lambda t^\alpha)
- E_\alpha(-\lambda s^\alpha)|
\leq C\,\lambda\,|t^\alpha - s^\alpha|$
uniformly in $\lambda$ (the $O(1/\lambda)$
contribution from the tail absorbs the prefactor at
infinity).  For $t, s \in [0, T]$:
\begin{itemize}[leftmargin=*]
\item If $\alpha \in (1/2, 1]$, subadditivity gives
  $|t^\alpha - s^\alpha| \leq |t-s|^\alpha$, so the
  deterministic Hölder exponent is
  $\beta_{\mathrm{det}} = \alpha$.
\item If $\alpha \in (1, 2)$, the mean value theorem
  gives $|t^\alpha - s^\alpha|
  \leq \alpha T^{\alpha-1}|t-s|$, so the
  deterministic part is Lipschitz
  ($\beta_{\mathrm{det}} = 1$).
\end{itemize}
Integrating against $|\hat{f}|^2$ and taking square
roots: $\norm{S_t f - S_s f}_{L^2}
\leq C\,|t-s|^{\beta_{\mathrm{det}}}\,\norm{f}_{L^2}$
with $\beta_{\mathrm{det}} = \min(\alpha, 1) \geq 1/2$
across the full parameter range.

\emph{Stochastic part.}  The factorization method
(as in \Cref{thm:continuity}) yields
$\beta_{\mathrm{stoch}} \leq 1/2 - \varepsilon$
under the integrability conditions of the applicable
existence theorem.

\emph{Combined.}  $\beta = \min(\beta_{\mathrm{det}},
\beta_{\mathrm{stoch}}) = \min(\alpha, 1/2-\varepsilon)
= \min(2H_2, 1/2-\varepsilon)$.  Since
$2H_2 > 1/2$ throughout the parameter range
$H_2 \in (1/4, 1)$, and $\beta_{\mathrm{det}} \geq 1/2
> 1/2 - \varepsilon$, the stochastic bound
determines $\beta$ everywhere; the $\min$ notation
is retained to emphasize that the deterministic
part could be the binding constraint in a regime
with smoother noise.
\end{proof}

\section{Energy, Decay, and Stationarity}\label{sec:energy}

\begin{theorem}[Energy identity]\label{thm:energy}
$E(t) := \E[\norm{X_t}_{L^2}^2] = \norm{S_t f}_{L^2}^2 + \int_0^t \norm{\Phi_s}_{\mathrm{HS}}^2\, ds$.
\end{theorem}

\begin{proof}
Expand using~\Cref{def:mild}; cross term vanishes ($\E[\delta(\cdot)] = 0$); It\^o isometry for the stochastic part.
\end{proof}

\begin{proposition}[Energy bound]\label{prop:dissipation}
For the deterministic equation $\partial_t^\alpha u = A_{H_1} u$, the energy $e(t) = \norm{u_t}_{L^2}^2$ satisfies $e(t) \leq e(0)$ for all $t \geq 0$, and is monotonically nonincreasing when $\alpha \in (0, 1]$.
\end{proposition}

\begin{remark}\label{rem:energy_inequality}
The identity $\partial_t^\alpha \norm{u}^2 = 2\ip{\partial_t^\alpha u}{u}$ fails for fractional derivatives.  The inequality $e(t) \leq e(0)$ follows directly from $|E_\alpha(-r)| \leq 1$ (\Cref{lem:ml_bounds}(i)), bypassing the product rule entirely.  For $\alpha \in (1, 2]$, $|E_\alpha(-r)|$ is not monotone in $r$ and $e(t)$ may oscillate within $[0, e(0)]$; the bound $e(t) \leq e(0)$ still holds.
\end{remark}

\begin{proof}
$e(t) = \int_\R |E_\alpha(-\lambda(\xi) t^\alpha)|^2 |\hat{f}(\xi)|^2\, d\xi \leq \int_\R |\hat{f}|^2\, d\xi = e(0)$, since $|E_\alpha(-r)| \leq 1$ by \Cref{lem:ml_bounds}(i).  For $\alpha \in (0, 1]$, $E_\alpha(-\lambda t^\alpha)$ is monotonically decreasing in $t$ (complete monotonicity), so $e(t)$ is monotonically nonincreasing.
\end{proof}

\begin{theorem}[Long-time decay]\label{thm:decay_det}
For $f \in L^1(\R) \cap L^2(\R)$ and $H_1 > 1/4$,
\begin{equation}\label{eq:decay}
\norm{S_t f}_{L^2} \leq C\,t^{-H_2/(2H_1)}
\norm{f}_{L^1} \qquad \text{as } t \to \infty.
\end{equation}
If additionally $\hat{f}(0) \neq 0$, the bound is
achieved: $\norm{S_t f}_{L^2} \geq c\,t^{-H_2/(2H_1)}$
for $t$ large.
\end{theorem}

\begin{proof}
\textbf{Upper bound.}  Write
$\norm{S_t f}_{L^2}^2
= \int_\R |E_\alpha(-\lambda(\xi) t^\alpha)|^2
|\hat{f}(\xi)|^2\, d\xi$.
Split at $\xi_* = t^{-\alpha/(2H_1)}$.

\emph{Low frequencies} ($|\xi| \leq \xi_*$):
$|E_\alpha(-r)| \leq 1$ and
$|\hat{f}(\xi)| \leq \norm{f}_{L^1}$, so the
contribution is $\leq \norm{f}_{L^1}^2 \cdot
2\xi_* = 2\norm{f}_{L^1}^2\,t^{-\alpha/(2H_1)}$.

\emph{High frequencies} ($|\xi| > \xi_*$):
$|E_\alpha(-r)| \leq C/r$ for $r \geq 1$
(Mittag-Leffler asymptotic), so
$|E_\alpha(-\lambda(\xi) t^\alpha)|^2
\leq C/(\lambda(\xi) t^\alpha)^2
= C/(c_{H_1}^2 |\xi|^{4H_1} t^{2\alpha})$.
The contribution is
$\leq C t^{-2\alpha} \int_{|\xi|>\xi_*}
|\xi|^{-4H_1} |\hat{f}(\xi)|^2\, d\xi$.
Split this region at $|\xi| = 1$.
On $\xi_* < |\xi| < 1$: use
$|\hat{f}(\xi)| \leq \norm{f}_{L^1}$ to get
$\leq t^{-2\alpha} \norm{f}_{L^1}^2
\int_{\xi_*}^1 \xi^{-4H_1}\, d\xi$.
For $H_1 > 1/4$,
$\int_{\xi_*}^1 \xi^{-4H_1}\, d\xi
\leq C\,\xi_*^{1-4H_1}
= C\,t^{-\alpha(1-4H_1)/(2H_1)}$.
Multiplying: $t^{-2\alpha} \cdot
t^{-\alpha(1-4H_1)/(2H_1)}
= t^{-\alpha/(2H_1)}$,
matching the low-frequency decay scale.
On $|\xi| \geq 1$:
$|\xi|^{-4H_1} \leq 1$, so the contribution is
$\leq t^{-2\alpha}\norm{f}_{L^2}^2$,
which is subdominant since
$2\alpha > \alpha/(2H_1)$ for $H_1 > 1/4$.

Combining: $\norm{S_t f}_{L^2}^2
\leq C\norm{f}_{L^1}^2\,t^{-\alpha/(2H_1)}
+ O(t^{-2\alpha})$.
For large $t$, the first term dominates since
$\alpha/(2H_1) < 2\alpha$ (equivalently $H_1 > 1/4$).
Taking the square root with $\alpha = 2H_2$:
$\norm{S_t f}_{L^2}
\leq C\,t^{-H_2/(2H_1)}\norm{f}_{L^1}$.

\textbf{Lower bound} (when $\hat{f}(0) \neq 0$).
By continuity of $\hat{f}$, there exists $\delta > 0$
with $|\hat{f}(\xi)| \geq |\hat{f}(0)|/2$ for
$|\xi| \leq \delta$.  Choose $r_0 > 0$ small enough
that $|E_\alpha(-r)| \geq 1/2$ for $r \in [0, r_0]$
(possible since $E_\alpha(0) = 1$ and $E_\alpha$ is
continuous; for $\alpha \in (0, 1]$ monotonicity
makes this immediate, and for $\alpha \in (1, 2]$
continuity at zero suffices).  Define
$\xi_{**} := \min(\delta,
(r_0/c_{H_1})^{1/(2H_1)}\,t^{-\alpha/(2H_1)})$.
For $t$ large enough that the second argument
is smaller than $\delta$,
$\xi_{**} = (r_0/c_{H_1})^{1/(2H_1)}\,
t^{-\alpha/(2H_1)}$, and on $|\xi| \leq \xi_{**}$
we have $\lambda(\xi) t^\alpha \leq r_0$, so
$|E_\alpha(-\lambda(\xi) t^\alpha)|^2 \geq 1/4$.
Therefore
\[
\norm{S_t f}_{L^2}^2
\geq \int_{|\xi| \leq \xi_{**}}
|E_\alpha(-\lambda(\xi) t^\alpha)|^2
|\hat{f}(\xi)|^2\, d\xi
\geq \frac{1}{4}\,(|\hat{f}(0)|/2)^2 \cdot 2\xi_{**}
= c\,t^{-\alpha/(2H_1)},
\]
which gives $\norm{S_t f}_{L^2}
\geq c\,t^{-H_2/(2H_1)}$ as claimed.
\end{proof}

\begin{remark}[Polynomial decay for all parameters]
The decay $t^{-H_2/(2H_1)}$ is \emph{always polynomial},
including the Markov case $H_2 = 1/2$, where the rate is
$t^{-1/(4H_1)}$.  On $\R$ without confining geometry,
there is no spectral gap and the spatial diffusion
produces algebraic (not exponential) decay even with
Markov time dynamics.  The distinction is that $H_2 < 1/2$
produces \emph{slower} polynomial decay than $H_2 = 1/2$,
reflecting the additional temporal memory.
\end{remark}

\begin{theorem}[Stationary covariance, conditional]
\label{thm:stationary}
The stationary covariance
\[
Q_\infty := \int_0^\infty \Phi_s\, Q\, \Phi_s^*\, ds
\]
exists as a trace-class operator on $L^2(\R)$ provided
\begin{equation}\label{eq:stationary_HS}
\int_0^\infty \norm{\Phi_s Q^{1/2}}_{\mathrm{HS}}^2\, ds
< \infty.
\end{equation}
For $Q$ a Fourier multiplier of nonnegative symbol
$q(\xi)$, this is equivalent (up to constants) to the
low-frequency integrability condition
\begin{equation}\label{eq:stationary_lowfreq}
\int_\R q(\xi)\, \lambda(\xi)^{-(2\alpha-1)/\alpha}\, d\xi
< \infty,
\qquad \lambda(\xi) = c_{H_1}|\xi|^{2H_1}.
\end{equation}
\end{theorem}

\begin{proof}
The Hilbert-Schmidt
condition~\eqref{eq:stationary_HS} is the Da~Prato--Zabczyk
trace-class criterion for the stochastic-convolution
covariance.  In Fourier-diagonal form,
\[
\sigma_\infty^2(\xi)
= q(\xi) \int_0^\infty s^{2(\alpha-1)}
|E_{\alpha,\alpha}(-\lambda(\xi) s^\alpha)|^2\, ds
= \frac{q(\xi)}{\alpha}\,
\lambda(\xi)^{-(2\alpha-1)/\alpha}
\int_0^\infty u^{1 - 1/\alpha}
|E_{\alpha,\alpha}(-u)|^2\, du,
\]
via the substitution $u = \lambda(\xi) s^\alpha$.  The
$u$-integral is finite for $\alpha > 1/2$ (integrable at
$0$ since $1 - 1/\alpha > -1$, and at infinity since
$|E_{\alpha,\alpha}(-u)|^2 = O(u^{-4})$ by the sharper
asymptotic of \Cref{lem:ml_bounds}(iii)).  Trace-class
$Q_\infty$ then reduces to
$\int \sigma_\infty^2(\xi)\, d\xi
< \infty$, which is~\eqref{eq:stationary_lowfreq}
up to constants.
\end{proof}

\begin{remark}[Stationary covariance is conditional on
low-frequency behavior]
\label{rem:stationary_conditional}
On $\R$, $\lambda(\xi) \to 0$ as $\xi \to 0$, so the
weight $\lambda(\xi)^{-(2\alpha-1)/\alpha}$ is
non-integrable at the origin.
Cylindrical noise ($q \equiv 1$) therefore fails
\eqref{eq:stationary_lowfreq} for every
$\alpha \in (1/2, 2)$, and a general trace-class $Q$
does not give a stationary covariance unless its symbol
$q(\xi)$ vanishes sufficiently fast at $\xi = 0$.  On a
bounded domain with positive first Dirichlet eigenvalue
$\lambda_1 > 0$ (so that $\lambda(\xi) \geq \lambda_1$
on the discrete spectrum), or under $q(\xi)$ with
sufficient low-frequency damping, the obstacle
disappears and $Q_\infty$ exists.
\end{remark}

\section{Variance Scaling}\label{sec:scaling}

\begin{proposition}[Variance scaling]\label{thm:selfsimilar}
Under the assumptions of \Cref{thm:existence_trace}
(trace-class noise) or \Cref{thm:existence}
(cylindrical noise with $H_1 > \alpha/(2(2\alpha-1))$),
for $f = 0$
and any $c > 0$, the variance scales as
\begin{equation}\label{eq:selfsimilar}
\E[\norm{X_{ct}}_{L^2}^2]
= c^{2H}\, \E[\norm{X_t}_{L^2}^2],
\qquad H = 2H_2 - \tfrac{1}{2} - \tfrac{H_2}{2H_1}.
\end{equation}
\end{proposition}

\begin{proof}
By the It\^o isometry (\Cref{cor:mode_variance}),
\[
\E[\norm{X_t}_{L^2}^2]
= \int_\R \int_0^t s^{2(\alpha-1)}
|E_{\alpha,\alpha}(-\lambda(\xi) s^\alpha)|^2\, ds\, d\xi.
\]
Substitute $s = t\sigma$:
$\E[\norm{X_t}_{L^2}^2]
= t^{2\alpha-1} \int_\R \int_0^1 \sigma^{2(\alpha-1)}
|E_{\alpha,\alpha}(-\lambda(\xi) t^\alpha
\sigma^\alpha)|^2\, d\sigma\, d\xi$.
Now substitute $\xi' = t^{\alpha/(2H_1)}\xi$, so
$\lambda(\xi) t^\alpha = \lambda(\xi')$ and
$d\xi = t^{-\alpha/(2H_1)}\,d\xi'$:
\[
\E[\norm{X_t}_{L^2}^2]
= t^{2\alpha - 1 - \alpha/(2H_1)}\,C,
\]
where $C = \int_\R \int_0^1 \sigma^{2(\alpha-1)}
|E_{\alpha,\alpha}(-\lambda(\xi')\sigma^\alpha)|^2\,
d\sigma\, d\xi'$ is independent of $t$.  The inner
integral converges since $2(\alpha-1) > -1$ (i.e.,
$H_2 > 1/4$) and $|E_{\alpha,\alpha}| \leq C/(1+r)$.
The outer integral converges by the same substitution
as in \Cref{thm:existence}.  Positivity $C > 0$ follows
from $E_{\alpha,\alpha}(0) = 1/\Gamma(\alpha) > 0$.  Replacing $t$ by $ct$:
the variance scales with exponent
$2H = 2\alpha - 1 - \alpha/(2H_1)
= 4H_2 - 1 - H_2/H_1$.
\end{proof}

\begin{remark}[Scaling exponent]
\label{rem:hurst}
The variance scaling exponent
$H = 2H_2 - 1/2 - H_2/(2H_1)$ depends on both
parameters.  In the Markov limit $H_2 = 1/2$:
$H = 1/2 - 1/(4H_1)$.  With the standard Laplacian
$H_1 = 1$: $H = 3H_2/2 - 1/2$.  The exponent
$H > 0$ (hence nontrivial scaling) requires
$H_2 > H_1/(4H_1 - 1)$ when $H_1 > 1/4$.
\end{remark}

\section{Special Cases}\label{sec:special}

Setting $H_2 = 1/2$ recovers Markov dynamics:
$\partial_t X = A_{H_1} X + \dot{B}$, with
polynomial decay $t^{-1/(4H_1)}$.
Setting $H_1 = 1$ gives the time-fractional heat
equation $\partial_t^\alpha X = -c\Delta X + \dot{B}$.
Both together ($H_1 = 1$, $H_2 = 1/2$) yield the
classical stochastic heat equation, whose
finite-dimensional projections are OU processes.
Restricting to a single Fourier mode $\xi_0$ gives
the scalar fractional OU process
$\partial_t^\alpha Y = -\lambda Y + \dot{W}$ with
$\lambda = c_{H_1}|\xi_0|^{2H_1}$,
studied by Cheridito, Kawaguchi, and
Maejima~\cite{cheridito2003}.

\section{The Operator Derivative on the Fractional Memory Space}
\label{sec:operator}

The operator derivative $D = \delta^*$, defined as the
adjoint of the stochastic integral
in~\cite{fontes2026opfac,fontes2026banach}, provides
representation theorems from a single duality identity.
Applied to the fractional memory process, it yields
exact Clark--Ocone representations in the Gaussian
case and identifies a structural defect under L\'evy
noise.  The interpretation of this defect as a
``memory gap'' is a modeling choice, not a theorem.

\subsection{Energy space and divergence}

The mild solution~(\Cref{def:mild}) writes
$X_t = S_t f + \delta(\Phi)$, where $\delta$ is the
stochastic integral (divergence) and $\Phi_{t-s}$ is
the Mittag-Leffler resolvent integrand.

\begin{definition}[Fractional memory energy space]
\label{def:energy_space}
The energy space at time $t$ is
\[
\HH_Q^t := L^2([0,t]; L^2(\R)),
\]
with divergence $\delta(u) := \int_0^t u_s\, dB_s$.
\end{definition}

\begin{definition}[Operator derivative]
\label{def:D_quijotic}
$D := \delta^*$, defined by the adjoint identity:
for $F \in L^2(\Omega)$ and $u \in \HH_Q^t$,
\begin{equation}\label{eq:adjoint_identity}
\ip{DF}{u}_{\HH_Q^t} = \E[F \cdot \delta(u)].
\end{equation}
\end{definition}

\begin{remark}[Volterra kernel structure]
The resolvent integrand $\widehat{\Phi_s}(\xi)
= s^{\alpha-1}
E_{\alpha,\alpha}(-\lambda(\xi) s^\alpha)$
is a Volterra kernel with power-law tail
$\sim s^{\alpha-1}/\Gamma(\alpha)$, reflecting the
Caputo time derivative.  This is the same kernel
structure encountered in the Volterra-L\'evy theory
of~\cite{fontes2026volterra}, with the Mittag-Leffler
function replacing the fractional power kernel.
\end{remark}

\subsection{Clark-Ocone in the Gaussian case}

\begin{theorem}[Exact representation]
\label{thm:clark_ocone}
In the Gaussian setting (noise $B$ is Brownian),
$D = \delta^*$ coincides with the Malliavin derivative,
the Leibniz rule holds (i.e., the Leibniz defect
$\Gamma(F,G) := D(FG) - F \cdot DG - G \cdot DF$
vanishes for all suitable $F, G$), and every
$F \in L^2(\FF_t)$ admits
\begin{equation}\label{eq:clark_ocone}
F - \E[F] = \int_0^t
\E[D_s F \mid \FF_s]\, dB_s.
\end{equation}
\end{theorem}

\begin{proof}
The Gaussian noise generates a Wiener space on which
$D = \delta^*$ is the Malliavin derivative
(see~\cite{nualart2006}, Chapter~1).  The Leibniz rule
follows from the product structure of the Gaussian
measure.  Clark-Ocone is the classical predictable
representation property on Wiener
space~\cite{nualart2006}, Theorem~1.3.14.
\end{proof}

\begin{remark}
In the Gaussian case, functionals of $X_t$ admit exact
stochastic integral representations against the input
stream.  The fractional memory process under Gaussian
noise is therefore the baseline of perfect
representability; the L\'evy case below shows where
this breaks down.
\end{remark}

\section{Extension to L\'evy Noise}
\label{sec:levy}

When the Gaussian noise $\dot{B}$ is replaced by a
L\'evy noise $\dot{L}$ with jumps, the Leibniz rule
fails and the defect acquires a concrete propagation
law governed by the Mittag-Leffler resolvent kernel.

Let $\dot{L}$ be a L\'evy noise with jump measure
$\nu$.  The equation becomes
\begin{equation}\label{eq:levy_fmp}
\partial_t^\alpha X_t = A_{H_1} X_t + \dot{L}_t,
\qquad X_0 = f.
\end{equation}
The mild solution retains the same resolvent structure:
\begin{equation}\label{eq:mild_levy}
X_t = S_t f + \int_0^t \Phi_{t-s}\, dL_s,
\end{equation}
where the stochastic integral is now against the
L\'evy process $L$.  Well-posedness of the mild
formulation under suitable moment conditions on $\nu$
follows from standard L\'evy stochastic integration
theory~\cite{daprato2014,meerschaert2012}.

\subsection{Combined energy space}

Following~\cite{fontes2026banach}, the energy space
splits into continuous and jump components:
\begin{equation}\label{eq:combined_space}
\HH = H_W \times \widetilde{H}_J,
\end{equation}
where $H_W = L^2([0,t]; L^2(\R))$ carries the Brownian
part and $\widetilde{H}_J$ carries the compensated
jump part, with norm weighted by the L\'evy measure
$\nu$ (see~\cite{fontes2026banach}, Definition~3.1, for
the precise construction).  The operator derivative
splits accordingly: $D = (D_W, \widetilde{D}_J)$.

\subsection{The Leibniz defect}

\begin{definition}[Leibniz defect]
\label{def:defect}
For $F, G \in \mathrm{dom}(D)$,
\begin{equation}\label{eq:defect}
\Gamma(F,G) := D(FG) - F \cdot DG - G \cdot DF.
\end{equation}
\end{definition}

\begin{theorem}[Defect dichotomy]
\label{thm:defect_dichotomy}
\leavevmode
\begin{enumerate}[label=(\roman*)]
\item If $\nu = 0$ (Gaussian noise):
  $\Gamma(F,G) = 0$ for all $F, G$.
\item If $\nu \neq 0$: $\Gamma$ is generically
  nonzero and is concentrated on the jump component
  $\widetilde{H}_J$.
\end{enumerate}
\end{theorem}

\begin{proof}
Part~(i) is \Cref{thm:clark_ocone}: the Gaussian
product structure ensures Leibniz holds for $D_W$.
Part~(ii) is proved in~\cite{fontes2026banach},
Theorem~4.2: the jump component $\widetilde{D}_J$
fails the Leibniz rule because the It\^o product
formula for jump processes introduces cross terms
between pre-jump and post-jump values that do not
factor through the derivative.  The defect $\Gamma$
is therefore supported on $\widetilde{H}_J$.
\end{proof}

\subsection{Defect propagation}

The resolvent structure of the fractional memory
process determines how the defect propagates in time.

\begin{proposition}[Defect propagation]
\label{prop:defect_propagation}
For the L\'evy fractional memory
process~\eqref{eq:levy_fmp}, a jump at time $s$
contributes to the defect at time $t > s$ with weight
\begin{equation}\label{eq:defect_weight}
|\widehat{\Phi_{t-s}}(\xi)|
\leq C\,(t-s)^{\alpha-1}
\cdot \frac{1}{1 + \lambda(\xi)(t-s)^\alpha}.
\end{equation}
For fixed $\xi$, this decays as $(t-s)^{\alpha-1}$
for small $t - s$ and as $(t-s)^{-1}$ for large
$t - s$---in both cases, polynomial.
\end{proposition}

\begin{proof}
Direct from the Mittag-Leffler bound
$|E_{\alpha,\alpha}(-r)| \leq C/(1+r)$ applied to the
resolvent integrand
$\widehat{\Phi_s}(\xi)
= s^{\alpha-1}
E_{\alpha,\alpha}(-\lambda(\xi) s^\alpha)$.
\end{proof}

\begin{remark}
In the Markov case ($\alpha = 1$), the exponential
kernel ensures rapid decay of the defect contribution.
For $\alpha < 1$, the Mittag-Leffler kernel decays
polynomially: the defect from a jump at time $s$
persists at time $t$ with weight
$\sim (t-s)^{-\alpha}$.
This parallels the unhedgeable risk in
Volterra-L\'evy markets~\cite{fontes2026volterra},
where the Leibniz defect equals the component of
financial risk that no trading strategy can eliminate.
\end{remark}

\begin{remark}
\label{rem:what_D_adds}
Compared to the standard Malliavin calculus, the
operator derivative framework adds the defect
$\Gamma$ as a structural invariant: its vanishing
characterizes the Gaussian case
(\Cref{thm:defect_dichotomy}), its $L^p$ norm
quantifies the deviation~\cite{fontes2026volterra},
and its propagation is governed by the resolvent
kernel~\eqref{eq:defect_weight}.
\end{remark}

\section{Modal Truncation and Discrete Representation}
\label{sec:discrete}

The spectral decomposition (\Cref{thm:spectral}) reduces
the infinite-dimensional SPDE to a continuum of scalar
modes.  In any finite implementation, only finitely many
modes are retained.  This section establishes the
effective temporal depth under modal truncation and
proves that, for a stream of discrete impulse inputs,
the spectrally truncated mild solution is computed
exactly by a finite Mittag-Leffler sum.

\subsection{Effective context under modal truncation}

\begin{definition}[Effective context]
\label{def:eff_context}
The effective context $C_{\mathrm{eff}}(N)$ is the
maximum temporal depth at which stored information
remains above threshold $\varepsilon > 0$:
$C_{\mathrm{eff}}(N) := \sup\{\tau > 0 :
\exists\, k \leq N \text{ with }
|E_\alpha(-\lambda_k \tau^\alpha)| > \varepsilon\}$,
where $\lambda_k = c_{H_1}|\xi_k|^{2H_1}$ for
$N$ uniformly spaced modes
$\xi_k = k\Xi/N$, $k = 1, \ldots, N$ (excluding the
zero mode, which has $\lambda(0) = 0$ and would not
decay).
\end{definition}

\begin{proposition}[Effective context scaling]
\label{thm:eff_context}
Under modal truncation to $N$ uniformly spaced
nonzero modes $\xi_k = k\Xi/N$, as $\varepsilon \to 0$,
\begin{equation}\label{eq:eff_context}
C_{\mathrm{eff}}(N) \lesssim
\left(\frac{N}{\Xi}\right)^{H_1/H_2}
\cdot \varepsilon^{-1/(2H_2)}.
\end{equation}
The bound is sharp (up to constants) when
$\alpha \in (0, 1]$, where
$E_\alpha(-\lambda \tau^\alpha)$ is monotonically
decreasing; for $\alpha \in (1, 2]$ the bound
controls the envelope of $|E_\alpha|$ but not the
precise crossing time, which is non-unique due to
oscillations.
\end{proposition}

\begin{proof}
The longest-lived mode (by envelope) is
$\xi_{\min} = \Xi/N$ with
$\lambda_{\min} = c_{H_1}(\Xi/N)^{2H_1}$.
By \Cref{lem:ml_bounds}(ii),
$|E_\alpha(-\lambda_{\min}\tau^\alpha)|
\leq C_\alpha/(1 + \lambda_{\min}\tau^\alpha)$.
This bound drops below $\varepsilon$ once
$\lambda_{\min}\tau^\alpha \gtrsim
\varepsilon^{-1}$, i.e.,
$\tau \gtrsim (\lambda_{\min})^{-1/\alpha}
\varepsilon^{-1/\alpha}
= (N/\Xi)^{H_1/H_2}\,\varepsilon^{-1/(2H_2)}$,
using $\alpha = 2H_2$.  Sharpness for $\alpha \in
(0, 1]$ follows from the matching asymptotic
$E_\alpha(-r) \sim 1/(r\Gamma(1-\alpha))$
(\cite[Ch.~4]{gorenflo2014}), which gives a lower
bound of the same order; for $\alpha \in (1, 2]$
the Mittag-Leffler function oscillates and there
are arbitrarily large $\tau$ with $|E_\alpha| <
\varepsilon$ interleaved with $|E_\alpha| \approx
\varepsilon$, so the crossing time is non-unique
but the envelope bound above still controls the
effective context in the sense of the
supremum in \Cref{def:eff_context}.
\end{proof}

\begin{remark}[Memory-compute tradeoff]
\label{rem:memory_compute}
The scaling $C_{\mathrm{eff}} \sim N^{H_1/H_2}$ is
polynomial in the number of retained modes.  Since
each mode evaluation costs $O(1)$ (the Mittag-Leffler
function can be evaluated via lookup table with
interpolation), the total query cost for $M$ stored
inputs is $O(NM)$.  The temporal depth thus grows as
a power law in the computational budget per query.
\end{remark}

\subsection{Discrete-impulse exactness}

\begin{theorem}[Discrete-impulse exactness]
\label{thm:discretization}
Consider the fractional memory
process~\eqref{eq:quijotic_main} with impulse forcing
at times $s_1 < \cdots < s_M < t$:
\[
\dot{B}_t = \sum_{j=1}^M v^{(j)}\,\delta(t - s_j),
\qquad v^{(j)} \in L^2(\R).
\]
Then the mild solution at time $t > s_M$ is
\begin{equation}\label{eq:mild_impulse}
X_t = S_t f + \sum_{j=1}^M \Phi_{t-s_j}\, v^{(j)},
\end{equation}
and its Fourier projection onto mode $\xi_k$ is
\begin{equation}\label{eq:discrete_exact}
\hat{X}_t(\xi_k)
= E_\alpha(-\lambda_k\,t^\alpha)\,\hat{f}(\xi_k)
+ \sum_{j=1}^M
(t - s_j)^{\alpha-1}\,
E_{\alpha,\alpha}\bigl(-\lambda_k\,(t - s_j)^\alpha\bigr)\,
\hat{v}^{(j)}(\xi_k),
\end{equation}
where $\lambda_k = c_{H_1}|\xi_k|^{2H_1}$.
\end{theorem}

\begin{proof}
Substitute $dB_s = \sum_j v^{(j)}\,\delta(s - s_j)\,ds$
into the mild-solution formula of \Cref{def:mild}.
The stochastic integral collapses by the sifting
property:
$\int_0^t \Phi_{t-s}\,dB_s
= \sum_j \Phi_{t-s_j}\,v^{(j)}$
for $t > s_M$, which gives~\eqref{eq:mild_impulse}.
Taking Fourier projection at mode $\xi_k$ and using
$\widehat{\Phi_\tau}(\xi) = \tau^{\alpha-1}
E_{\alpha,\alpha}(-\lambda(\xi)\tau^\alpha)$
(\Cref{def:resolvent_kernel})
and $\widehat{S_t f}(\xi_k)
= E_\alpha(-\lambda_k t^\alpha)\,\hat{f}(\xi_k)$
(\Cref{def:resolvent})
yields~\eqref{eq:discrete_exact}.
\end{proof}

\begin{remark}[Two representations of impulse
response]
\label{rem:two_representations}
Formula~\eqref{eq:discrete_exact} is the exact mild
solution of the forced
equation~\eqref{eq:quijotic_main} and uses the
two-parameter function $E_{\alpha,\alpha}$ with
prefactor $(t-s_j)^{\alpha-1}$.  This is the forced
Green's function of the fractional SPDE.  For
$\alpha \in (1/2, 1)$ (i.e., $H_2 \in (1/4, 1/2)$),
the prefactor is singular as $t \to s_j^+$,
reflecting the weak regularity of the impulse
response at the time of impulse.

A distinct representation, natural for memory
architectures that treat each stored input as an
initial condition placed on the free dynamics at
the time of writing, is the
\emph{autonomous-decay representation}:
\begin{equation}\label{eq:autonomous_decay}
\hat{X}^{\mathrm{auto}}_t(\xi_k)
= \sum_{j=1}^M
E_\alpha\bigl(-\lambda_k\,(t - s_j)^\alpha\bigr)\,
\hat{v}^{(j)}(\xi_k),
\end{equation}
which uses the one-parameter $E_\alpha$ and has no
prefactor.  Formula~\eqref{eq:autonomous_decay} is
\emph{not} the mild solution
of~\eqref{eq:quijotic_main}: it corresponds to the
jump-reset model in which, at each $s_j$, the state
is augmented by $v^{(j)}$ and then evolves
autonomously under
$\partial_t^\alpha X = A_{H_1}X$ until the next
impulse.  The two representations coincide when
$\alpha = 1$ (both give $e^{-\lambda_k(t-s_j)}$) and
disagree for $\alpha \neq 1$.  The companion
paper~\cite{fontes2026arch} uses
\eqref{eq:autonomous_decay} as the retrieval kernel
on the grounds that it matches the operation
``each stored fact contributes a decaying
relevance weight to a ranking at query time'';
we defer to that paper for discussion of the
modeling choice, which is a choice distinct from the
forced-SPDE impulse response derived here.
\end{remark}

\begin{corollary}[$L^1$ integrability separates the
two representations]
\label{cor:l1_separation}
Fix $\lambda > 0$ and consider the scalar kernels on
$[0, \infty)$:
\[
G_{\mathrm{forced}}(\tau) := \tau^{\alpha-1}
E_{\alpha,\alpha}(-\lambda \tau^\alpha),
\qquad
G_{\mathrm{auto}}(\tau) := E_\alpha(-\lambda \tau^\alpha).
\]
\begin{enumerate}[label=\textnormal{(\roman*)}]
\item For $\alpha \in (0, 1]$,
  $G_{\mathrm{forced}} \geq 0$ and
  $\int_0^\infty G_{\mathrm{forced}}(\tau)\,d\tau
  = 1/\lambda$.  For $\alpha \in (1, 2)$,
  $\|G_{\mathrm{forced}}\|_{L^1([0,\infty))}
  < \infty$.
\item For $\alpha \in (0, 1)$,
  $G_{\mathrm{auto}} \notin L^1([0, \infty))$: the
  integral $\int_0^\infty E_\alpha(-\lambda \tau^\alpha)
  \,d\tau$ diverges at infinity.  For $\alpha = 1$,
  $G_{\mathrm{auto}} = e^{-\lambda \tau}$ with
  $L^1$ norm $1/\lambda$.  For $\alpha \in (1, 2)$,
  $G_{\mathrm{auto}} \in L^1([0, \infty))$.
\end{enumerate}
The $L^1$ integrability transition of
$G_{\mathrm{auto}}$ occurs at $\alpha = 1$, coinciding
with the Markov-fractional boundary.
\end{corollary}

\begin{proof}
\emph{(i), case $\alpha \in (0, 1]$.}  Complete
monotonicity of $E_\alpha(-t^\alpha)$ implies
$E_{\alpha,\alpha}(-r) \geq 0$ for $r \geq 0$
(via the identity $t^{\alpha-1} E_{\alpha,\alpha}
(-\lambda t^\alpha) = -\lambda^{-1}\partial_t
E_\alpha(-\lambda t^\alpha)$ applied to the
nonincreasing $E_\alpha$), so
$G_{\mathrm{forced}} \geq 0$.  By monotone
convergence,
$\int_0^\infty G_{\mathrm{forced}}(\tau)\,d\tau
= \lim_{s \to 0^+} \mathcal{L}[G_{\mathrm{forced}}](s)
= \lim_{s \to 0^+} 1/(s^\alpha + \lambda)
= 1/\lambda$.

\emph{(i), case $\alpha \in (1, 2)$.}
$G_{\mathrm{forced}}$ is continuous on $(0, \infty)$
and $\tau^{\alpha-1}$ is integrable near zero
(since $\alpha > 0$).  At infinity,
\Cref{lem:ml_bounds}(iii) gives
$|E_{\alpha,\alpha}(-\lambda \tau^\alpha)|
\leq C/(1 + \lambda \tau^\alpha)$, so
$|G_{\mathrm{forced}}(\tau)|
\leq C \tau^{\alpha-1}/(1 + \lambda \tau^\alpha)
= O(\tau^{-1})$ at infinity.  The sharper
asymptotic $E_{\alpha,\alpha}(-r) = O(r^{-2})$
(remark after \Cref{lem:ml_bounds}) improves this
to $O(\tau^{-\alpha-1})$, which is integrable for
$\alpha > 0$.

\emph{(ii), case $\alpha \in (0, 1)$.}  By the
Mittag-Leffler asymptotic,
$E_\alpha(-\lambda \tau^\alpha)
\sim 1/[\lambda\Gamma(1-\alpha)\,\tau^\alpha]$ as
$\tau \to \infty$, which is positive for $\alpha
\in (0, 1)$ (since $\Gamma(1-\alpha) > 0$ there).
$\int^\infty \tau^{-\alpha}\,d\tau$ diverges iff
$\alpha \leq 1$.  Combined with $E_\alpha(0) = 1$,
this gives $\int_0^\infty G_{\mathrm{auto}}\,d\tau
= +\infty$.

\emph{(ii), case $\alpha = 1$.}  Direct.

\emph{(ii), case $\alpha \in (1, 2)$.}  The same
asymptotic gives $|E_\alpha(-\lambda\tau^\alpha)|
\sim 1/[|\lambda\Gamma(1-\alpha)|\,\tau^\alpha]$ with
$\alpha > 1$ (so $\tau^{-\alpha}$ is integrable at
infinity).  Bounded on compact sets (continuity).
Hence $L^1$.
\end{proof}

\begin{remark}[Rankings, aggregates, and the
complementary pathologies]
\label{rem:l1_consequence}
Corollary~\ref{cor:l1_separation} pins down when the
two representations of \Cref{rem:two_representations}
behave differently, and when they do not.

For an operation that is a function of ratios of
finitely many kernel evaluations --- top-$k$
retrieval ranking, pairwise comparisons, relative
ordering --- the $L^1$ norm of the kernel does not
appear directly, but the ranking still depends on
the kernel's pointwise shape.  Two kernels with
the same $L^1$ norm can produce different rankings;
two kernels with different $L^1$ norms can produce
the same rankings.  The $L^1$ question is separate
from the ranking question.

For an operation that is a function of
\emph{aggregate kernel mass} --- normalized mixture
weights $G(t - s_j)/\sum_i G(t - s_i)$ as the number
of stored facts $M \to \infty$, prior integrals
$\int_0^\infty G(\tau)\,\pi(\tau)\,d\tau$ under
heavy-tailed priors on age, time-integrated
expectations --- the $L^1$ norm appears and the
representations diverge qualitatively.  In the
polynomial-retention regime $\alpha \in (0, 1)$ the
forced-SPDE kernel is $L^1$-normalizable (total
mass $1/\lambda$) while the autonomous-decay kernel
is not.  Under a model in which writes arrive as a
stationary point process with density $\rho > 0$,
the aggregate weight
$\sum_j G_{\mathrm{auto}}(t - s_j)$ diverges as the
observation window grows, while the analogous sum
under $G_{\mathrm{forced}}$ remains bounded.

However, the forced-SPDE kernel carries its own
pathology: the $(t - s)^{\alpha-1}$ prefactor is
singular as $t \to s^+$ for $\alpha \in (0, 1)$
(\Cref{rem:two_representations}).  In a retrieval
system where queries may arrive close in time to
write events, this singularity causes
$G_{\mathrm{forced}}$ to rank the most recent write
first regardless of the semantic content of the
query, overwhelming any pointwise semantic
similarity signal by a factor that grows
unboundedly as the query-write gap shrinks.

The two kernels thus play \emph{complementary}
rather than hierarchical roles.  $G_{\mathrm{forced}}$
has $L^1$-bounded aggregate mass but is
pathological at zero age; it is suitable for
aggregate operations when queries are separated
from writes by a timescale on the order of
$\tau_{\mathrm{char}} = \lambda^{-1/\alpha}$, and
unsuitable when they are not.
$G_{\mathrm{auto}}$ is well-behaved at zero age but
$L^1$-divergent at infinity for $\alpha \in (0, 1)$;
it is suitable for direct retrieval ranking but
requires finite-horizon truncation or acceptance of
unboundedness for aggregate operations.

A cleaner statement of this complementarity uses
$L^p$ regularity.  For $\alpha \in (0, 1)$, direct
computation from the Mittag-Leffler asymptotics
$E_\alpha(-\lambda\tau^\alpha) \sim
1/(\lambda\Gamma(1-\alpha)\tau^\alpha)$ and
$E_{\alpha,\alpha}(-\lambda\tau^\alpha) \sim
\alpha/(\Gamma(1-\alpha)(\lambda\tau^\alpha)^2)$
at infinity, together with the boundary behaviour
$E_\alpha(0) = 1$ and
$E_{\alpha,\alpha}(0) = 1/\Gamma(\alpha)$, yields
\begin{equation}
G_{\mathrm{auto}} \in L^p([0,\infty))
\iff p > \tfrac{1}{\alpha},
\qquad
G_{\mathrm{forced}} \in L^p([0,\infty))
\iff p < \tfrac{1}{1-\alpha}.
\end{equation}
The critical exponents $1/\alpha$ and $1/(1-\alpha)$
are H\"older conjugates, since $\alpha + (1-\alpha) =
1$.  The two kernels thus live on opposite sides
of H\"older-conjugate endpoints: $G_{\mathrm{auto}}$
is in $L^\infty$ but not $L^1$,
$G_{\mathrm{forced}}$ is in $L^1$ but not $L^\infty$,
and for $\alpha \in (1/2, 1)$ they share the
open window $(1/\alpha, 1/(1-\alpha)) \ni 2$ while
for $\alpha \in (0, 1/2)$ the two windows are
disjoint.

This duality aligns with a natural distinction
between retrieval operators.  Sup-type operators
--- $\operatorname{argmax}$, top-$k$, threshold
gates, streaming $\limsup$ --- require only that
the kernel be bounded on its domain, i.e., $L^\infty$,
and $G_{\mathrm{auto}}$ provides this for all
$\alpha \in (0, 1)$ while $G_{\mathrm{forced}}$
does not.  Sum-type operators --- softmax
normalization, attention mixtures, Bayesian
posteriors over age, softmax-weighted RAG contexts
--- require the kernel to be summable against the
density of stored facts, which reduces to $L^1$
regularity under a stationary write process; and
$G_{\mathrm{forced}}$ provides this for all
$\alpha \in (0, 1)$ while $G_{\mathrm{auto}}$
does not.  The choice of temporal kernel and the
choice of retrieval operator are linked, not
independent: a pipeline restricted to sup-type
aggregation lives in $L^\infty$ and selects
$G_{\mathrm{auto}}$, and the $L^1$ divergence of
Corollary~\ref{cor:l1_separation} never arises;
a pipeline using sum-type aggregation lives in
$L^1$ and selects $G_{\mathrm{forced}}$ modulo its
zero-age singularity.

Applications requiring both bounded aggregates and
regular zero-age behaviour simultaneously are not
served by either kernel and may motivate a third
choice outside the Mittag-Leffler family --- for
example, a stretched exponential
$\exp\bigl(-(\lambda\tau)^\alpha\bigr)$, which is
$L^1$-integrable for $\alpha \in (0, 1)$ and finite
at $\tau = 0$ but lacks an SPDE derivation; or a
regularized forced-SPDE kernel
$(\tau + \epsilon)^{\alpha-1}
 E_{\alpha,\alpha}(-\lambda(\tau + \epsilon)^\alpha)$
with $\epsilon > 0$, which is bounded at $\tau = 0$
while retaining $L^1$-integrability.  The present
paper does not develop either alternative.
\end{remark}

The next two subsections turn the foregoing
asymptotic statements into checkable numerical
content: \Cref{sec:numerical_verification} converts
the $L^1$ separation into a closed-form
aggregate-mass identity under a stationary Poisson
write process, and \Cref{sec:recency_pathology}
exhibits the recency-singularity pathology of
$G_{\mathrm{forced}}$ in a synthetic retrieval
pipeline.

\subsection{Numerical verification of the
\texorpdfstring{$L^1$}{L1} separation}
\label{sec:numerical_verification}

The qualitative and quantitative content of
\Cref{cor:l1_separation,rem:l1_consequence} can be
illustrated by the closed-form aggregate-mass
expectations under a stationary Poisson write process
on $[0, T]$ with rate $\rho$ and decay $\lambda$.
By Campbell's theorem,
\begin{align}
\E[S_{\mathrm{forced}}(T)]
&= \rho \int_0^T (T-s)^{\alpha-1}
   E_{\alpha,\alpha}(-\lambda(T-s)^\alpha)\,ds
 = \frac{\rho}{\lambda}
   \bigl(1 - E_\alpha(-\lambda T^\alpha)\bigr),
\label{eq:s_forced_expectation}\\
\E[S_{\mathrm{auto}}(T)]
&= \rho \int_0^T E_\alpha(-\lambda u^\alpha)\,du
 = \rho\, T\, E_{\alpha,2}(-\lambda T^\alpha),
\label{eq:s_auto_expectation}
\end{align}
where the first identity is the antiderivative
relation
\begin{equation}\label{eq:E_alpha_derivative}
\frac{d}{dT} E_\alpha(-\lambda T^\alpha)
= -\lambda\,T^{\alpha-1}\,
E_{\alpha,\alpha}(-\lambda T^\alpha),
\end{equation}
which integrated on $[0, T]$ with
$E_\alpha(0) = 1$ gives
$\int_0^T u^{\alpha-1}
E_{\alpha,\alpha}(-\lambda u^\alpha)\,du
= [1 - E_\alpha(-\lambda T^\alpha)]/\lambda$;
the second follows from the Mittag-Leffler integral
identity $\int_0^T E_\alpha(-\lambda u^\alpha)\,du
= T\,E_{\alpha,2}(-\lambda T^\alpha)$~\cite[\S 4]{gorenflo2014}.
The total $L^1$ mass of $G_{\mathrm{forced}}$ on
$[0, \infty)$ equals $1/\lambda$
(\Cref{cor:l1_separation}(i)) and is recovered by
sending $T \to \infty$ in
\eqref{eq:s_forced_expectation}, equivalently by
evaluating $\mathcal{L}[\tau^{\alpha-1}
E_{\alpha,\alpha}(-\lambda \tau^\alpha)](s)
= 1/(s^\alpha + \lambda)$ at $s \to 0^+$.
The aggregate-mass ratio is therefore
\begin{equation}\label{eq:l1_ratio}
\mathcal{R}_\alpha(T)
:= \frac{\E[S_{\mathrm{auto}}(T)]}
        {\E[S_{\mathrm{forced}}(T)]}
 = \frac{\lambda\,T\,E_{\alpha,2}(-\lambda T^\alpha)}
        {1 - E_\alpha(-\lambda T^\alpha)},
\end{equation}
and the Mittag-Leffler asymptotics
$E_{\alpha,2}(-r)\sim 1/[\Gamma(2-\alpha)\,r]$ and
$E_\alpha(-r)\sim 1/[\Gamma(1-\alpha)\,r]$ as
$r \to \infty$ for $\alpha \in (0,1)$ give
\begin{equation}\label{eq:l1_ratio_asymptotic}
\mathcal{R}_\alpha(T)
\sim \frac{T^{1-\alpha}}{\Gamma(2-\alpha)}
\qquad \text{as } T \to \infty.
\end{equation}

\begin{figure}[h]
\centering
\includegraphics[width=\textwidth]{figures/fig_l1_mass}
\caption{\textbf{Closed-form aggregate kernel mass
under a Poisson write process,
$\rho = 1$, $\lambda = 0.05$.}
Blue:
$\E[S_{\mathrm{auto}}](T) = \rho T\,
E_{\alpha,2}(-\lambda T^\alpha)$
(\eqref{eq:s_auto_expectation}).
Red:
$\E[S_{\mathrm{forced}}](T) = \rho(1 -
E_\alpha(-\lambda T^\alpha))/\lambda$
(\eqref{eq:s_forced_expectation}).
Dotted horizontal line: $\rho/\lambda = 20$.
Dotted gray: reference $T^{1-\alpha}$ slope.
Across $\alpha \in \{0.3, 0.5, 0.8\}$ the forced-SPDE
mass saturates at $\rho/\lambda$ as predicted by
\Cref{cor:l1_separation}(i), while the autonomous-decay
mass diverges at rate $T^{1-\alpha}$ per
\eqref{eq:l1_ratio_asymptotic}.  Curves are deterministic
and exactly reproducible from
\texttt{scripts/generate\_fig\_l1\_mass.py}.}
\label{fig:l1_mass}
\end{figure}

The slope of $\E[S_{\mathrm{auto}}](T)$ on log--log
axes over the rightmost decade $T \in [10^7, 10^8]$,
computed from \eqref{eq:s_auto_expectation}, is
$\{0.732, 0.502, 0.200\}$ for $\alpha \in \{0.3, 0.5, 0.8\}$
respectively, against the asymptotic prediction
$\{0.700, 0.500, 0.200\}$; the small residual at
$\alpha = 0.3$ is the sub-asymptotic correction in the
Mittag-Leffler expansion at finite $T$.  The closed-form
ratio $\mathcal{R}_\alpha(T)$ at $T = 10^8$ takes the
values $\{4.30 \cdot 10^5,\ 1.13 \cdot 10^4,\ 43.4\}$
for the same $\alpha$, agreeing with the asymptotic
form \eqref{eq:l1_ratio_asymptotic} to within $2\%$.
The full table of slopes and ratios is written by the
generation script to
\texttt{outputs/l1\_mass\_metrics.txt}.

The practical consequence is as predicted by
\Cref{rem:l1_consequence}: at $T = 10^8$ and
$\alpha = 0.3$, $\mathcal{R}_\alpha(T) \approx 4.3 \cdot 10^5$,
and the two representations continue to separate at
rate $T^{1-\alpha}$.  Any operation depending on the
ratio $w(\tau)/S(T)$ --- normalized attention mixtures,
Bayesian age posteriors, softmax-style calibration
--- will therefore yield asymptotically different
answers under the two representations.  Rank-based
retrieval at fixed store size is insensitive to the
$L^1$ divergence per se, but is sensitive to the
pointwise shape of the kernel, in particular to the
recency singularity of $G_{\mathrm{forced}}$, as we
show next.

\subsection{Verification of the recency-singularity
pathology}
\label{sec:recency_pathology}

The penultimate paragraph of \Cref{rem:l1_consequence}
asserts that $G_{\mathrm{forced}}$'s
$(t-s)^{\alpha-1}$ prefactor causes it to dominate
over semantic similarity signals in a retrieval
context when queries arrive close to writes.  We
verify this claim on a synthetic retrieval pipeline.

The setup: $N$ documents at Poisson-arrival times on
$[0, T]$ with embeddings drawn uniformly from the
unit sphere in $\R^{d_{\mathrm{emb}}}$ ($d_{\mathrm{emb}}
= 64$).  Each query embedding is a target document's
embedding plus Gaussian noise at level $0.35$,
renormalized.  Retrieval score for each document is
the product of cosine similarity (shifted to be
nonnegative) and the temporal weight under either
$G_{\mathrm{auto}}$ or $G_{\mathrm{forced}}$.  Two
conditions: \emph{no lag} (query at time $T$,
youngest document at age near $0$) and \emph{lag}
(query at time $T + 10$, youngest document at age
$\geq 10$).

\begin{figure}[h]
\centering
\includegraphics[width=\textwidth]{figures/fig_retrieval_pipeline}
\caption{\textbf{Retrieval pipeline with semantic
$\times$ temporal scoring, $\alpha = 0.5$,
$\lambda = 0.05$.}  Top row: query at time $T$
(query-write gap can be near zero).  Bottom row:
query at time $T + 10$.  Left column: top-$1$
accuracy.  Middle: median softmax confidence $p_1$.
Right: $N_{\mathrm{eff}}$.  Under forced-SPDE with
no lag, top-$1$ accuracy is essentially zero across
all store sizes because the singular prefactor
$(t-s)^{\alpha-1}$ assigns near-infinite weight to
the youngest document regardless of semantic match.
With lag $=10$ the pathology softens but
$G_{\mathrm{forced}}$ still underperforms
$G_{\mathrm{auto}}$ on ranking by a factor of $4$ or
more.  The $L^1$-related predictions (bounded
$N_{\mathrm{eff}}$ under $G_{\mathrm{forced}}$,
linear growth under $G_{\mathrm{auto}}$) are
confirmed in the right column of both rows.}
\label{fig:pipeline}
\end{figure}

The experiment confirms the architectural picture of
\Cref{rem:l1_consequence}: the two kernels have
complementary pathologies that make them
non-interchangeable.  The $L^1$-based predictions of
\Cref{cor:l1_separation} hold under both conditions
in the right-column panels, but the left column
shows that $G_{\mathrm{forced}}$'s bounded
$N_{\mathrm{eff}}$ comes at the cost of misplaced
confidence: its high $p_1$ in the middle column is
confidence in the youngest document rather than the
semantically matching one.  A practitioner seeking
an $L^1$-integrable, non-recency-pathological
temporal weighting must look outside the
Mittag-Leffler family, as noted at the end of
\Cref{rem:l1_consequence}.

\begin{remark}
As $N \to \infty$ (more modes) and $M \to \infty$
(denser write events), the discrete
representation~\eqref{eq:discrete_exact} recovers
the continuum mild solution in $L^2(\R)$.  The
convergence rate is $O(\xi_{\min}^{-2H_1})$ in the
temporal direction (set by the slowest retained
mode) and depends on the sampling density $M/T$
in the input direction.
\end{remark}

\section{Conclusion}\label{sec:conclusion}

We have developed the fractional memory process---the
mild solution of $\partial_t^{2H_2} X_t = A_{H_1} X_t
+ \dot{B}_t$---as a mathematically rigorous model of
persistent-memory dynamics with independent spatial and
temporal parameters.

Well-posedness holds for $H_2 > 1/4$ under trace-class
noise and additionally $H_1 > \alpha/(2(2\alpha-1))$
(forcing $H_2 > 1/3$) under cylindrical
noise, with pathwise continuity established via the
factorization method.  The spectral decomposition
reduces the SPDE to independent scalar fractional OU
modes, yielding sharp decay $t^{-H_2/(2H_1)}$ (with
matching lower bound) and variance scaling exponent
$H = 2H_2 - 1/2 - H_2/(2H_1)$.  On the operator-theoretic side, $D = \delta^*$ yields exact
Clark--Ocone representations under Gaussian noise
($\Gamma = 0$) and a nonzero, polynomially propagating
Leibniz defect under L\'evy noise ($\Gamma \neq 0$).
Finally, the discrete-impulse exactness theorem shows
that for impulse inputs, the spectrally truncated mild
solution is a finite sum of two-parameter
Mittag-Leffler terms at $O(NM)$ cost, with effective
temporal depth scaling as $N^{H_1/H_2}$.  The
memory-architecture kernel used
in~\cite{fontes2026arch} is a distinct
autonomous-decay representation (\Cref{rem:two_representations})
that uses the one-parameter Mittag-Leffler function
and coincides with the forced-SPDE response only at
$\alpha = 1$.  The two representations differ
qualitatively in $L^1$ integrability
(\Cref{cor:l1_separation}): the forced-SPDE kernel
has total mass $1/\lambda$ across all $\alpha \in
(0, 2]$, while the autonomous-decay kernel is not
$L^1$-integrable for $\alpha \in (0, 1)$.  Under a
stationary Poisson write process, this produces a
power-law $T^{1-\alpha}$ separation in total kernel
mass between the two representations, verified
numerically in \Cref{sec:numerical_verification}.

The companion paper~\cite{fontes2026arch} develops
a concrete memory architecture using the
autonomous-decay kernel.  The forced-SPDE kernel
derived here is not a drop-in replacement for
retrieval: its $(t-s)^{\alpha-1}$ prefactor is
singular as $t \to s^+$ for $\alpha \in (0, 1)$,
causing it to dominate over semantic similarity
signals when queries arrive close to writes
(\Cref{sec:recency_pathology}).  The two kernels
thus play complementary rather than hierarchical
roles --- $G_{\mathrm{forced}}$ is $L^1$-bounded
but pathological at zero age, while
$G_{\mathrm{auto}}$ is regular at zero age but
$L^1$-divergent for $\alpha \in (0, 1)$.
Applications requiring both bounded aggregate mass
and regular zero-age behaviour may motivate a third
choice outside the Mittag-Leffler family, a
direction we leave to future work.

\subsection*{Acknowledgments}

AI tools were used during drafting and mathematical
exploration; the author takes full responsibility for
all content and errors.
\begin{thebibliography}{20}

\bibitem{fontes2026arch}
R.~Fontes, Fractional memory for attention-based
systems: architecture and experiments,
Quijotic Research, 2026.

\bibitem{fontes2026opfac}
R.~Fontes, The operator derivative in continuous
stochastic calculus: a Hilbert energy space framework,
Quijotic Research, 2026.
arXiv:2601.09976; Zenodo DOI:10.5281/zenodo.14618498.

\bibitem{fontes2026banach}
R.~Fontes, Operator factorization beyond Hilbert spaces:
representability obstructions and Leibniz defects for
stable L\'evy processes, Quijotic Research, 2026.
Zenodo DOI:10.5281/zenodo.14798498.

\bibitem{fontes2026volterra}
R.~Fontes, The boundary of hedgeability: pricing and
hedging in Volterra-L\'evy markets, Quijotic Research,
2026.
Zenodo DOI:10.5281/zenodo.19211991.

\bibitem{mainardi2010}
F.~Mainardi, \emph{Fractional Calculus and Waves in
Linear Viscoelasticity}, Imperial College Press, 2010.

\bibitem{podlubny1999}
I.~Podlubny, \emph{Fractional Differential Equations},
Academic Press, 1999.

\bibitem{gorenflo2014}
R.~Gorenflo et al., \emph{Mittag-Leffler Functions,
Related Topics and Applications}, Springer, 2014.

\bibitem{meerschaert2012}
M.~M.~Meerschaert and A.~Sikorskii, \emph{Stochastic
Models for Fractional Calculus}, De Gruyter, 2012.

\bibitem{daprato2014}
G.~Da~Prato and J.~Zabczyk, \emph{Stochastic Equations
in Infinite Dimensions}, 2nd ed., Cambridge, 2014.

\bibitem{pollard1948}
H.~Pollard, The completely monotonic character of
$E_\alpha(-x)$, \emph{Bull.\ AMS}, 54:1115--1116, 1948.

\bibitem{schneider1996}
W.~R.~Schneider, Completely monotone generalized
Mittag-Leffler functions,
\emph{Expo.\ Math.}, 14:3--16, 1996.

\bibitem{cheridito2003}
P.~Cheridito, H.~Kawaguchi, and M.~Maejima, Fractional
Ornstein-Uhlenbeck processes, \emph{EJP}, 8:1--14, 2003.

\bibitem{nualart2006}
D.~Nualart, \emph{The Malliavin Calculus and Related
Topics}, 2nd ed., Springer, 2006.

\end{thebibliography}

\end{document}
